\chapter*{Conclusion Générale}
\addcontentsline{toc}{chapter}{Conclusion Générale}
\markboth{Conclusion Générale}{Conclusion Générale}

\section*{Synthèse de la Démarche}

Ce projet a cherché à répondre à une problématique concrète : comment permettre aux
chercheurs et aux étudiants des institutions académiques tunisiennes de provisionner
et gérer de façon autonome des machines virtuelles et des clusters Kubernetes sur
l'infrastructure Huawei Cloud Stack (HCS) du MESRS, sans dépendre d'une intervention
manuelle d'un administrateur ?

La réponse apportée est \textbf{VM Marketplace}, une plateforme cloud full-stack en
libre-service développée sur six mois de stage chez Huawei Technologies Tunisie.
En travaillant de manière itérative à travers trois sprints, nous avons bâti une
architecture à trois services --- un frontend Next.js, un backend FastAPI et un moteur
d'inférence IA distinct --- qui intègre le paiement Stripe, le provisionnement
d'infrastructure IaC via Terraform sur HCS, la diffusion SSE pour le suivi en temps
réel, des recommandations de VM assistées par IA via un modèle Llama~3.1 hébergé
localement, un accès SSH dans le navigateur, la supervision Netdata et un gestionnaire
complet du cycle de vie des clusters Kubernetes.

\section*{Principaux Résultats}

\begin{itemize}
  \item Le provisionnement bout en bout de VM sur l'environnement HCS de production
        a été validé, avec des VM atteignant l'état \texttt{Running} et la supervision
        Netdata installée automatiquement via SSH.
  \item Un cluster Kubernetes à deux nœuds (K8s v1.29.15) a été provisionné avec
        succès, les deux nœuds atteignant l'état \texttt{Ready}. L'architecture réseau
        maître-comme-passerelle-NAT --- combinant \texttt{source\_dest\_check=false}
        au niveau Terraform avec des règles \texttt{iptables} MASQUERADE Linux ---
        a résolu la contrainte de quota EIP sans nécessiter de ressources cloud
        supplémentaires.
  \item Le moteur de recommandation IA, fonctionnant entièrement sur une instance
        Ollama locale, a correctement mappé des descriptions textuelles de charges de
        travail vers des recommandations de flavor HCS avec score de confiance.
  \item Les cinq objectifs initiaux ont été atteints et validés sur l'infrastructure
        de production.
\end{itemize}

\section*{Difficultés et Enseignements}

Le défi technique le plus significatif de ce projet a été la découverte et la
résolution du comportement \texttt{source\_dest\_check} au niveau de l'hyperviseur
HCS. Cette contrainte n'est pas documentée dans la documentation standard du fournisseur
Terraform \texttt{huaweicloud/hcs} et n'a été identifiée qu'après une élimination
systématique des autres causes possibles (règles iptables, paramètres noyau, DNS).
L'enseignement retenu est que les comportements spécifiques à l'hyperviseur cloud
requièrent une investigation empirique et ne peuvent pas être supposés conformes aux
comportements génériques de la mise en réseau Linux.

Sur le plan du génie logiciel, ce projet a renforcé la valeur d'une approche itérative
lorsqu'on travaille dans un environnement à inconnues externes : chaque sprint a validé
un ensemble d'hypothèses sur l'environnement HCS et a adapté la conception en
conséquence.

\section*{Compétences Acquises}

Ce projet m'a permis de construire et de consolider les compétences suivantes :

\begin{itemize}
  \item \textbf{Infrastructure cloud :} gestion des ressources Huawei Cloud Stack,
        Terraform IaC, configuration réseau VPC, gestion des EIP et configuration
        cloud-init.
  \item \textbf{Administration Kubernetes :} initialisation de cluster avec kubeadm,
        déploiement du CNI Flannel, opérations kubectl et gestion du kubeconfig.
  \item \textbf{Développement backend :} API REST FastAPI, Server-Sent Events,
        authentification JWT, gestion de threads en arrière-plan, automatisation SSH
        Paramiko.
  \item \textbf{Intégration IA :} ingénierie de prompts LLM avec sortie JSON
        structurée, hébergement de modèle local avec Ollama, systèmes de recommandation
        à score de confiance.
  \item \textbf{Intégration full-stack :} coordination de trois services indépendants
        (frontend, backend, moteur IA) avec un contrat d'API partagé.
\end{itemize}

\section*{Apport pour l'Organisme d'Accueil}

VM Marketplace fournit à Huawei Technologies Tunisie un prototype fonctionnel de portail
cloud en libre-service adapté aux spécificités de l'environnement HCS MESRS. La
solution réseau pour les nœuds workers Kubernetes privés documentée dans ce rapport
(\texttt{source\_dest\_check} + iptables MASQUERADE) constitue une connaissance
technique réutilisable pour les futurs déploiements HCS. La plateforme démontre
également qu'une couche de recommandation IA basée sur un LLM hébergé localement est
viable sans dépendre d'API IA cloud externes, ce qui est particulièrement pertinent
dans un contexte de données sensibles de recherche.

\section*{Perspectives et Travaux Futurs}

Quatre améliorations concrètes sont proposées pour une évolution vers une version de
production :

\begin{enumerate}
  \item \textbf{Persistance du suivi de progression :} remplacer les dictionnaires
        \texttt{progress\_store} en mémoire par un magasin Redis ou une table journal
        en base de données afin de survivre aux redémarrages du backend pendant les
        opérations de provisionnement longues.
  \item \textbf{Suite de tests automatisés :} implémenter des tests unitaires pour
        la couche service (en simulant les appels de sous-processus Terraform et les
        connexions SSH) et des tests d'intégration contre un projet HCS de staging.
  \item \textbf{RBAC multi-tenant :} ajouter des rôles utilisateur (administrateur,
        utilisateur standard) et des quotas de ressources par utilisateur appliqués
        au niveau de la plateforme, indépendamment des quotas du projet HCS.
  \item \textbf{Migration PostgreSQL et containerisation :} remplacer SQLite par
        PostgreSQL et containeriser les trois services avec Docker Compose ou un chart
        Helm Kubernetes, permettant la mise à l'échelle horizontale et le déploiement
        multi-hôtes.
\end{enumerate}

\section*{Réflexion Finale}

Construire une plateforme qui réduit l'écart entre un environnement cloud de production
et des utilisateurs académiques non experts a nécessité de naviguer en permanence entre
la conception de l'expérience utilisateur, l'ingénierie des systèmes distribués et
les contraintes techniques spécifiques de l'infrastructure HCS. Ce stage a démontré
que combler ces couches --- rendre les ressources cloud puissantes accessibles aux
chercheurs qui souhaitent simplement exécuter leur code --- est à la fois
techniquement exigeant et véritablement utile.
